% Source: https://artofproblemsolving.com/community/c3h2168972_2020_mock_amc_11_problem_21_discussion
% Stable identity: aops-topic-2168972-q21
Let Alice and Bob run $a$ and $b$ laps per second. Running in opposite directions, they pass each other exactly when the total distance they have covered, $(a+b)t$ laps, is a whole number. So
\[(a+b)\cdot\dfrac{70}{27}=k\quad\text{for some positive integer }k,\qquad a+b=\dfrac{27k}{70}.\]
The lap counts give $15\le60a<16$ and $8\le60b<9$, so $23\le60(a+b)<25$. Now $60(a+b)=\dfrac{162k}{7}$, which is about $23.14$ for $k=1$ and more than $46$ for $k\ge2$. Hence $k=1$ and $a+b=\dfrac{27}{70}$.

To make $a$ as large as possible, make $b$ as small as possible: $b\ge\dfrac{8}{60}=\dfrac{2}{15}$, so
\[a\le\dfrac{27}{70}-\dfrac{2}{15}=\dfrac{81-28}{210}=\dfrac{53}{210}.\]
This value is possible: with $a=\dfrac{53}{210}$ and $b=\dfrac{2}{15}$, Alice runs $60a=\dfrac{106}{7}\approx15.14$ laps and Bob exactly $8$ laps in $60$ seconds, and they pass each other at $t=\dfrac{70}{27}$. Since $53$ is prime and does not divide $210$, the fraction is in lowest terms, and $m+n=53+210=\boxed{263}$, answer $\textbf{(C)}$.
